%% =========================================================
%% Physics example: the weighted Ising bond-object twist as an
%% instance of the multi-block asymmetric compression theorem.
%% Covers Notes/OpenProblemsTN/checks/asym_ising_action_data.md,
%% verified by checks/asym_ising_action_verify.py, and the P6
%% renormalization-fixed-point twist of
%% Notes/OpenProblemsTN/problems/p6_rfp_structure_constant_l_dependence.tex.
%% =========================================================

\section{The Ising anyon chain: exact arithmetic in \texorpdfstring{$\Z[\sqrt2]$}{Z[sqrt(2)]}}%
\label{sec:asymex_ising_ring}

The topological-symmetry tensors of the Ising anyon chain have entries built
from the F-symbols of the Ising fusion category. The quantum dimension of
the non-abelian object is $\sqrt2$, and after rescaling, every entry lies
in $\Z[\sqrt2]$. The embedding of this ring into $\C$ preserves the
matrix-unit spanning criterion of Theorem~\ref{thm:asymex_normal_units}.
The tensors are introduced in Section~\ref{sec:asymex_ising_tensors}.
Here the spanning criterion establishes injectivity at one site. It does
not impose the spectral-radius-one normalization required of a normal
tensor (NT) in the CPSV convention.

\begin{definition}[The ring of \texorpdfstring{$\sqrt2$}{sqrt(2)}-integers and its complex embedding]%
    \label{def:asymex_zsqrt2_ring}
    \lean{zsqrt2ToComplex, MPSTensor.complexOfZsqrt2, MPSTensor.mulZsqrt2Tensor,
        MPSTensor.actZsqrt2Tensor}
    \leanok
    The ring $\Z[\sqrt2]$ consists of the numbers $a+b\sqrt2$ with
    $a,b\in\Z$, with multiplication determined by $(\sqrt2)^2=2$.
    It embeds into $\C$ by sending $\sqrt2$ to the positive real square
    root of two. Applying this embedding entrywise gives a complex matrix
    from a matrix over $\Z[\sqrt2]$. For matrix product operator tensors
    $M$ and $N$ over $\Z[\sqrt2]$, their local product is
    $(M\mathbin{\cdot}N)^{ik}=\sum_jM^{ij}\otimes N^{jk}$.
    For a matrix product operator tensor $M$ and a state tensor $A$,
    the local action is $(M\cdot A)^i=\sum_jM^{ij}\otimes A^j$.
\end{definition}

\begin{theorem}[Exact arithmetic in \texorpdfstring{$\Z[\sqrt2]$}{Z[sqrt(2)]}]%
    \label{thm:asymex_zsqrt2_ring}
    % Principal declaration: MPSTensor.complexOfZsqrt2_mul.
    \lean{zsqrt2ToComplex_injective, zsqrt2ToComplex_ne_zero,
        MPSTensor.complexOfZsqrt2_mul, MPSTensor.mulTensor_complexOfZsqrt2,
        MPSTensor.actTensor_complexOfZsqrt2,
        MPSTensor.isNormal_of_single_eq_smul_zsqrt2}
    \leanok
    \uses{def:asymex_zsqrt2_ring}
    The embedding of $\Z[\sqrt2]$ into $\C$ is injective and sends every
    nonzero element to a nonzero complex number. The entrywise embedding
    of matrices is multiplicative and commutes with the local product
    and local action of tensors over $\Z[\sqrt2]$.
    If the one-site matrices of a tensor over $\Z[\sqrt2]$ include a
    nonzero $\Z[\sqrt2]$-multiple of every matrix unit, its complex
    embedding is injective at one site.
\end{theorem}

\begin{proof}
    \leanok
    Injectivity follows from the uniqueness of the integer coordinates
    $a,b$ representing $a+b\sqrt2$: if $a+b\sqrt2=a'+b'\sqrt2$ as real
    numbers, then $(b-b')\sqrt2=a'-a$. Irrationality of $\sqrt2$ gives
    $b=b'$, whence $a=a'$. Multiplicativity of the entrywise
    embedding follows because it is induced by a ring homomorphism.
    The two compatibility statements follow by expanding the local product
    and local action entrywise. The spanning argument is the same as in
    Theorem~\ref{thm:asymex_normal_units}: dividing each designated
    one-site matrix by its nonzero scalar factor exhibits every matrix
    unit in the span.
\end{proof}

\section{The three topological-symmetry tensors of the Ising anyon chain}%
\label{sec:asymex_ising_tensors}

The Ising fusion category has three simple objects, the vacuum $1$, the
fermion $\psi$, and the non-abelian anyon $\sigma$, with fusion rules
$\psi\otimes\psi=1$, $\sigma\otimes\psi=\sigma$ and
$\sigma\otimes\sigma=1\oplus\psi$. Each object $a$ gives a matrix product
operator tensor $A_a$ on the ten fusion-tree labels $h=(x',\rho,x)$ with
$x\in x'\otimes\rho$. Its local matrices are determined by the F-symbols
$[F^{a\,x'\,\rho}_y]_{y'x}$; a local matrix of $A_a$ vanishes unless
its two physical labels carry the same middle label $\rho$.

\begin{definition}[The three tensors \texorpdfstring{$A_1$, $A_\psi$, $A_\sigma$}{A1, Apsi, Asigma}]%
    \label{def:asymex_ising_tensors}
    \lean{IsingTwist.isingRho, IsingTwist.isingOneZ, IsingTwist.isingPsiZ,
        IsingTwist.isingSigmaZ, IsingTwist.isingOne, IsingTwist.isingPsi,
        IsingTwist.isingSigma}
    \leanok
    \uses{def:asymex_zsqrt2_ring, def:mpo_tensor}
    The ten fusion-tree labels, indexed from $0$ to $9$ in order, are
    $(1,1,1)$, $(1,\psi,\psi)$, $(1,\sigma,\sigma)$, $(\psi,1,\psi)$,
    $(\psi,\psi,1)$, $(\psi,\sigma,\sigma)$, $(\sigma,1,\sigma)$,
    $(\sigma,\psi,\sigma)$, $(\sigma,\sigma,1)$ and $(\sigma,\sigma,\psi)$.
    Their middle labels are
    $1,\psi,\sigma,1,\psi,\sigma,1,\psi,\sigma,\sigma$.
    The tensor $A_1$ has auxiliary space $P_1$ with ordered basis
    $[(1,1),(\psi,\psi),(\sigma,\sigma)]$
    and has a matrix unit at each of its ten diagonal physical-index pairs.
    The tensor $A_\psi$ has auxiliary space $P_\psi$ with ordered basis
    $[(1,\psi),(\psi,1),(\sigma,\sigma)]$
    and has a signed matrix unit at each of ten physical-index pairs.
    The sign at
    $((\sigma,\psi,\sigma),(\sigma,\psi,\sigma))$
    is the F-symbol $[F^{\psi\sigma\psi}_\sigma]=-1$.
    The rescaled tensor $\sqrt2A_\sigma$ has auxiliary space $P_\sigma$
    with ordered basis
    $[(1,\sigma),(\psi,\sigma),(\sigma,1),(\sigma,\psi)]$.
    It has sixteen nonzero local matrices, each a matrix unit scaled by
    $\pm1$ or $\pm\sqrt2$, from the rescaled Ising F-symbols.
    In particular, the factors $\pm1$ arise from the rescaled
    $\sigma\sigma\sigma$ F-symbols $\pm1/\sqrt2$.
\end{definition}

\begin{theorem}[Sector vanishing of the three tensors]%
    \label{thm:asymex_ising_sector_vanishing}
    \lean{IsingTwist.isingOneZ_eq_zero_of_rho_ne,
        IsingTwist.isingPsiZ_eq_zero_of_rho_ne,
        IsingTwist.isingSigmaZ_eq_zero_of_rho_ne}
    \leanok
    \uses{def:asymex_ising_tensors}
    For each of the three tensors, the local matrix at a pair of labels
    with distinct middle labels vanishes.
\end{theorem}

\begin{proof}
    \leanok
    Inspecting the ten or sixteen nonzero local matrices in
    Definition~\ref{def:asymex_ising_tensors} shows that each is indexed
    by labels with the same middle label.
\end{proof}

\begin{theorem}[One-site injectivity of the three tensors]%
    \label{thm:asymex_ising_normal}
    \lean{IsingTwist.isingOne_isNormal, IsingTwist.isingPsi_isNormal,
        IsingTwist.isingSigma_isNormal}
    \leanok
    \uses{def:asymex_ising_tensors, def:normal}
    Each of the three tensors is injective at one site: the ten nonzero
    local matrices of $A_1$ span $\MN{3}$, those of $A_\psi$ span
    $\MN{3}$, and the sixteen nonzero local matrices of $\sqrt2A_\sigma$
    span $\MN{4}$.
\end{theorem}

\begin{proof}
    \leanok
    \uses{thm:asymex_zsqrt2_ring}
    Every matrix unit of the relevant auxiliary matrix algebra occurs,
    up to a nonzero factor in $\Z[\sqrt2]$, among the local matrices.
    This is a finite verification over integer or $\Z[\sqrt2]$-valued
    entries. Theorem~\ref{thm:asymex_zsqrt2_ring} gives the corresponding
    spanning statement over $\C$.
\end{proof}

\section{The weighted direct sum and its local product}%
\label{sec:asymex_ising_bond_object}

Consider the weighted direct sum
$\Theta=\lambda_1A_1\oplus\lambda_\psi A_\psi\oplus\lambda_\sigma A_\sigma$
and its product with $A_\sigma$. We first use only the two summands
$\lambda_1A_1\oplus\lambda_\psi A_\psi$, at
$(\lambda_1,\lambda_\psi)=(1/2,3/10)$, rescaled to the integer weights
$5,3$. Denote this rescaled direct sum by $\Theta_2$.
Rescaling $A_\sigma$ by $\sqrt2$ puts every entry in
$\Z[\sqrt2]$.

The tensor $\Theta_2\star(\sqrt2A_\sigma)$ represents a product of two
matrix product operators on the pair alphabet of
Section~\ref{sec:asym_mpo_product}, not an operator applied to a state.
Here $\star$ denotes the local product defined above.
The relevant result is Theorem~\ref{thm:asym_mpo_product}.
Section~\ref{sec:asymex_ising_action} treats the separate action of the
full weighted direct sum on a fusion-tree product state, using
Corollary~\ref{cor:asym_action_tensors}.
These explicit tensor identities do not by themselves establish a
positive vertical decomposition of an MPDO renormalization fixed point.

\begin{definition}[The weighted direct sum and the product tensor]%
    \label{def:asymex_ising_bond_object}
    \lean{IsingTwist.thetaTwoZ, IsingTwist.thetaTwo, IsingTwist.isingBondObjectZ,
        IsingTwist.isingBondObject}
    \leanok
    \uses{def:asymex_ising_tensors, def:mpdo_operator_product}
    Set $\Theta_2=5A_1\oplus3A_\psi$, with auxiliary space
    $\C^6=P_1\oplus P_\psi$. The local product tensor
    $B=\Theta_2\star(\sqrt2A_\sigma)$ has a hundred-letter pair alphabet
    and auxiliary dimension $24=6\cdot4$, with coordinate order
    $4\,(\Theta_2\text{ index})+(P_\sigma\text{ index})$.
\end{definition}

\begin{center}
    % Formula: O_L(Theta_2) O_L(sqrt2 A_sigma) = O_L(B), the stacked product
    %   of Definition~\ref{def:asymex_ising_bond_object} closed into full
    %   periodic operators (data file, section 1.2, and Theorem~\ref{thm:asymex_ising_trace}
    %   below).
    % Ink-to-index map: the upper row of operator nodes is Theta_2, the
    %   lower row is sqrt2 A_sigma; the physical legs h_1,...,h_L of the top
    %   row and g_1,...,g_L of the bottom row are open, and the shared
    %   virtual wire between the two rows at every site is the summed middle
    %   label of the stacked product. On the right, the single row of B
    %   carries both physical labels at each site as one pair letter.
    % Contraction set: the L vertical wires joining the two rows on the
    %   left, one per site; the horizontal virtual wires within each row,
    %   closed periodically.
    % Boundary signature (virtual-west, virtual-east, physical-up,
    %   physical-down): left, (0,0,L,L) with each three-site representative
    %   emitting (0,0,3,3); right, (0,0,L,0) with the representative emitting
    %   (0,0,3,0), the pair label absorbing both physical legs into one node.
    \tenkzkernel
    \begin{tenkz}[rows={op,op}, cols=4, boundary=periodic]
        \tn[skin=mpo, ports={90:physical:$h_1$, 270:physical}]{\Theta_2} &
        \tn[skin=mpo, ports={90:physical:$h_2$, 270:physical}]{\Theta_2} &
        \tn[skin=dots, ports={180:virtual, 0:virtual}]{} &
        \tn[skin=mpo, ports={90:physical:$h_L$, 270:physical}]{\Theta_2} \\
        \tn[skin=mpo, ports={90:physical, 270:physical:$g_1$}]{\sqrt2A_\sigma} &
        \tn[skin=mpo, ports={90:physical, 270:physical:$g_2$}]{\sqrt2A_\sigma} &
        \tn[skin=dots, ports={180:virtual, 0:virtual}]{} &
        \tn[skin=mpo, ports={90:physical, 270:physical:$g_L$}]{\sqrt2A_\sigma}
    \end{tenkz}
    \quad $=$ \quad
    \begin{tenkz}[rows={op}, cols=4, boundary=periodic]
        \tn[skin=mpo, ports={90:physical:$h_1$, 270:physical:$g_1$}]{B} &
        \tn[skin=mpo, ports={90:physical:$h_2$, 270:physical:$g_2$}]{B} &
        \tn[skin=dots, ports={180:virtual, 0:virtual}]{} &
        \tn[skin=mpo, ports={90:physical:$h_L$, 270:physical:$g_L$}]{B}
    \end{tenkz}
\end{center}

\section{The gauge and the block decomposition}%
\label{sec:asymex_ising_gauge}

The change of auxiliary coordinates is a signed permutation of the
twenty-four coordinates of $B$. Its first four rows are the F-move
isometry for $1\otimes\sigma=\sigma$, the next four are the F-move
isometry for $\psi\otimes\sigma=\sigma$, and the remaining sixteen
form a coordinate complement. The single negative sign is the F-symbol
$[F^{\psi\sigma\psi}_\sigma]=-1$.

\begin{definition}[Signed permutation matrices and the Ising gauge]%
    \label{def:asymex_ising_gauge}
    \lean{MPSTensor.signedPermMatrix, IsingTwist.isingPerm, IsingTwist.isingSign,
        IsingTwist.isingGaugeZ, IsingTwist.isingGauge, IsingTwist.isingSlots,
        IsingTwist.isingWeights, IsingTwist.isingTargets, IsingTwist.isingBlockZ}
    \leanok
    \uses{def:asymex_ising_bond_object, def:asymex_explicit_gauge}
    For a map $\pi$ on an index set and a sign vector $s$, form the
    matrix with entry $s(x)$ in row $x$ and column $\pi(x)$ and zero
    elsewhere. When $\pi$ is a permutation, this is a signed permutation
    matrix. The Ising gauge is the signed permutation matrix given by the
    F-move rows described above, in the sense of
    Definition~\ref{def:asymex_explicit_gauge}.
    There are two nonzero blocks: one from $1\otimes\sigma=\sigma$,
    with weight $5$, and one from $\psi\otimes\sigma=\sigma$, with
    weight $3$. Both carry $\sqrt2A_\sigma$, of auxiliary dimension
    four. The zero block has dimension $16=24-2\cdot4$.
\end{definition}

\begin{theorem}[Orthogonality of the Ising gauge]%
    \label{thm:asymex_ising_gauge_unitary}
    \lean{MPSTensor.signedPermMatrix_mul_mul_transpose,
        IsingTwist.isingGaugeZ_mul_transpose, IsingTwist.isingGaugeZ_transpose_mul}
    \leanok
    \uses{def:asymex_ising_gauge}
    For a permutation $\pi$ and sign vector $s$, let
    $G_{x,\pi(x)}=s(x)$ and let all other entries of $G$ vanish.
    For every matrix $M$,
    \begin{align}
        (GMG^{\mathsf T})_{xy} & =s(x)\,M_{\pi(x)\pi(y)}\,s(y).
        \label{eq:asymex_ising_gauge_conj}
    \end{align}
    The Ising gauge $G$ satisfies $GG^{\mathsf T}=G^{\mathsf T}G=\Id$,
    so its transpose is its inverse.
\end{theorem}

\begin{proof}
    \leanok
    Identity~(\ref{eq:asymex_ising_gauge_conj}) follows by expanding the
    matrix product against the two Kronecker deltas of the signed
    permutation matrices. For the Ising gauge, both products with its
    transpose are the identity, by finite verification over the
    twenty-four rows.
\end{proof}

\begin{theorem}[The local block identity for the Ising product tensor]%
    \label{thm:asymex_ising_letter_identity}
    % Principal declaration: IsingTwist.isingGaugeZ_conj.
    \lean{IsingTwist.isingGaugeZ_conj, IsingTwist.isingBondObject_conjMatrix}
    \leanok
    \uses{def:asymex_ising_gauge}
    In the block coordinates of Definition~\ref{def:asymex_ising_gauge},
    conjugation by the Ising gauge gives, for every pair $(h,h')$,
    \begin{align}
        GB^{(h,h')}G^{-1}
        &=5(\sqrt2A_\sigma)^{(h,h')}
          \oplus3(\sqrt2A_\sigma)^{(h,h')}\oplus0_{16}.
    \end{align}
\end{theorem}

\begin{proof}
    \leanok
    \uses{thm:asymex_ising_sector_vanishing, thm:asymex_ising_gauge_unitary}
    A local matrix indexed by labels with distinct middle labels vanishes by
    Theorem~\ref{thm:asymex_ising_sector_vanishing}, as does the corresponding
    block on the right. For labels sharing a middle label, the conjugation
    formula of Theorem~\ref{thm:asymex_ising_gauge_unitary} expresses the
    conjugated matrix in permuted coordinates. The local product restricts
    to a sum over labels in the shared sector. The identity then follows
    by finite verification of the entries in
    Definitions~\ref{def:asymex_ising_tensors}
    and~\ref{def:asymex_ising_gauge}, treating the three middle-label
    sectors separately.
\end{proof}

\section{Compression onto the \texorpdfstring{$\sigma$}{sigma} tensor}%
\label{sec:asymex_ising_compression}

\begin{theorem}[Compression of the Ising product tensor]%
    \label{thm:asymex_ising_compression}
    % Principal declaration: IsingTwist.isingCompression.
    \lean{IsingTwist.isingCompression,
        IsingTwist.isingBondObject_dim_eq, IsingTwist.isingBondObject_z_eq}
    \leanok
    \uses{thm:asymex_ising_letter_identity, thm:asym_mpo_product,
        thm:asym_multiblock}
    The tensor $B$ of Definition~\ref{def:asymex_ising_bond_object}
    has the block triangular form of Theorem~\ref{thm:asym_multiblock}
    with two weighted copies of $\sqrt2A_\sigma$, of weights $5$ and
    $3$, and a sixteen-dimensional zero block, $24=4+4+16$.
    The compression pairs are mutually biorthogonal and satisfy
    $W_sB^{(h,h')}V_s=(\mu_s\sqrt2A_\sigma)^{(h,h')}$
    for every pair of physical labels, where $\mu_0=5$ and $\mu_1=3$.
\end{theorem}

\begin{proof}
    \leanok
    \uses{thm:asymex_explicit_gauge, thm:asymex_zsqrt2_ring}
    Theorem~\ref{thm:asymex_ising_letter_identity} gives the block-diagonal
    form of every conjugated local matrix over $\Z[\sqrt2]$.
    Theorem~\ref{thm:asymex_zsqrt2_ring} preserves this identity over
    $\C$. Block diagonality gives triangularity and identifies the two
    nonzero diagonal blocks with the weighted copies of $\sqrt2A_\sigma$,
    by Theorem~\ref{thm:asymex_explicit_gauge}.
    The dimension count is clause (vii) of
    Theorem~\ref{thm:asym_multiblock}.
\end{proof}

\begin{theorem}[The extension splits]%
    \label{thm:asymex_ising_split}
    % Principal declaration: IsingTwist.isingBondObject_remainder.
    \lean{IsingTwist.isingBondObject_remainder}
    \leanok
    \uses{thm:asymex_ising_compression, prop:asym_projector_weighted_sum}
    The remainder of the compression in
    Theorem~\ref{thm:asymex_ising_compression} vanishes identically.
    Thus the extension splits, and each nonzero block has local
    intertwiners in both directions. For every pair $(h,h')$,
    \begin{align}
        B^{(h,h')}V_s&=V_s(\mu_s\sqrt2A_\sigma)^{(h,h')},\\
        W_sB^{(h,h')}&=(\mu_s\sqrt2A_\sigma)^{(h,h')}W_s.
    \end{align}
\end{theorem}

\begin{proof}
    \leanok
    \uses{thm:asymex_ising_letter_identity}
    Every conjugated local matrix is block diagonal by
    Theorem~\ref{thm:asymex_ising_letter_identity}. Its off-diagonal
    blocks, which form the conjugated remainder, therefore vanish.
    Restricting the block-diagonal identity to each nonzero block gives
    the two intertwining identities.
\end{proof}

\begin{center}
    % Formula: B V_s = V_s C_s and W_s B = C_s W_s of
    %   Theorem~\ref{thm:asymex_ising_split}, with C_s = mu_s sqrt2 A_sigma.
    % Ink-to-index map: the ring node V_s (resp. W_s) is the compression map
    %   into (resp. out of) the bond space; the labelled node B carries the
    %   physical letter h and the open virtual legs are the bond spaces of B
    %   on the left and of the slot on the right.
    % Contraction set: the single virtual wire joining B and the ring node on
    %   each side; no physical contraction.
    % Boundary signature (virtual-west, virtual-east, physical-up,
    %   physical-down): each equality is (1,1,1,0).
    \tenkzkernel
    \begin{tenkz}[rows={wire}, cols=2, west=open, east=open]
        \tn[label pos=s, ports={90:physical:$h$}]{B} & \tn[skin=ring]{V_s}
    \end{tenkz}
    \quad $=$ \quad
    \begin{tenkz}[rows={wire}, cols=2, west=open, east=open]
        \tn[skin=ring]{V_s} & \tn[label pos=s, ports={90:physical:$h$}]{C_s}
    \end{tenkz}
    \\[2ex]
    \begin{tenkz}[rows={wire}, cols=2, west=open, east=open]
        \tn[skin=ring]{W_s} & \tn[label pos=s, ports={90:physical:$h$}]{B}
    \end{tenkz}
    \quad $=$ \quad
    \begin{tenkz}[rows={wire}, cols=2, west=open, east=open]
        \tn[label pos=s, ports={90:physical:$h$}]{C_s} & \tn[skin=ring]{W_s}
    \end{tenkz}
\end{center}

\section{The length-dependent product coefficient}%
\label{sec:asymex_ising_trace}

\begin{theorem}[The trace identity for the Ising product tensor]%
    \label{thm:asymex_ising_trace}
    \lean{IsingTwist.isingBondObject_trace_evalWord, IsingTwist.isingTwist_mpo}
    \leanok
    \uses{thm:asymex_ising_compression, prop:asym_compression_trace}
    For every nonempty word $w$ over the alphabet of pair labels,
    \begin{align}
        \tr(B^w) & =(5^{|w|}+3^{|w|})\tr((\sqrt2A_\sigma)^w).
        \label{eq:asymex_ising_trace}
    \end{align}
    Equivalently, at every positive length $L$, the periodic operators
    satisfy $O_L(\Theta_2)O_L(\sqrt2A_\sigma)=(5^L+3^L)O_L(\sqrt2A_\sigma)$.
\end{theorem}

\begin{proof}
    \leanok
    Proposition~\ref{prop:asym_compression_trace} applied to the
    compression in Theorem~\ref{thm:asymex_ising_compression} gives
    $\tr(B^w)$ as the sum of the two weighted block traces.
    Pulling each weight out of the matrix product contributes
    $5^{|w|}$ and $3^{|w|}$. Reading this identity entrywise at the
    pairs of physical configurations gives the operator identity.
\end{proof}

\begin{remark}[Two fusion channels give the same output tensor]%
    \label{rem:asymex_ising_hidden_channel}
    The fusions $1\otimes\sigma=\sigma$ and
    $\psi\otimes\sigma=\sigma$ give two copies of the same output tensor
    with different weights. The coefficient in
    (\ref{eq:asymex_ising_trace}) is the power sum
    $5^L+3^L=\tr(\diag(5,3)^L)$.
    This is a length-dependent coefficient for the stated periodic
    product identity. Its power-sum form alone does not identify these
    tensors with a basis of normal tensors (BNT) in the positive vertical
    decomposition of an MPDO renormalization fixed point, nor prove that
    the length dependence cannot be removed under other choices of tensors.
\end{remark}

\begin{center}
    % Formula: the physics reading of Remark~\ref{rem:asymex_ising_hidden_channel},
    %   the strand sqrt2 A_sigma with the weighted bond object
    %   Theta = 5 Pi_1 + 3 Pi_psi inserted at one bond.
    % Ink-to-index map: each labelled node is sqrt2 A_sigma, physical legs
    %   omitted; the ring node on the inner bond is the weighted sum of the
    %   two rank-one idempotents cut out by the two hidden bond objects.
    % Contraction set: the horizontal virtual wires joining consecutive
    %   sqrt2 A_sigma nodes, closed periodically through the ring node.
    % Boundary signature: (virtual-west=0, virtual-east=0, physical-up=3,
    %   physical-down=0) for this three-site representative of the periodic
    %   chain; the ellipsis carries the remaining sites.
    \tenkzkernel
    \begin{tenkz}[cols=6, boundary=periodic, physical=up]
        \tn[label pos=135, ports={90:physical:$h_1$}]{\sqrt2A_\sigma} &
        \tn[skin=ring]{\Theta} &
        \tn[label pos=135, ports={90:physical:$h_2$}]{\sqrt2A_\sigma} &
        \tn[skin=ring]{\Theta} &
        \tn[skin=dots]{} & \tn[skin=ring]{\Theta}
    \end{tenkz}
    \qquad $\Theta=5\,\Pi_1+3\,\Pi_\psi$
\end{center}

\section{Action on a fusion-tree product state}%
\label{sec:asymex_ising_action}

We now apply the full weighted direct sum
$\Theta_3=5A_1\oplus3A_\psi\oplus2A_\sigma$ to a fusion-tree product
state, using Corollary~\ref{cor:asym_action_tensors}.
This is an operator acting on a state, rather than the operator product
of Section~\ref{sec:asymex_ising_bond_object}.

\begin{definition}[The full weighted direct sum and the fusion-tree product states]%
    \label{def:asymex_ising_action}
    \lean{IsingTwist.thetaThreeZ, IsingTwist.thetaThree, IsingTwist.productStateZ,
        IsingTwist.productState, IsingTwist.sectorState}
    \leanok
    \uses{def:asymex_ising_tensors, def:mps_tensor}
    Set $\Theta_3=5A_1\oplus3A_\psi\oplus2A_\sigma$, with auxiliary space
    $\C^{10}=P_1\oplus P_\psi\oplus P_\sigma$.
    For a fusion-tree label $h_0$, let $\psi_{h_0}$ be the
    bond-dimension-one tensor generating $\ket{h_0}^{\otimes L}$
    at every length. In particular,
    $\ket{\sigma1\sigma}^{\otimes L}$ is the product state at the
    label $(\sigma,1,\sigma)$, the fusion tree of the trivial anyon chain.
\end{definition}

\begin{theorem}[One-site injectivity of the fusion-tree product-state tensors]%
    \label{thm:asymex_ising_action_normal}
    \lean{IsingTwist.productState_isNormal}
    \leanok
    \uses{def:asymex_ising_action, def:normal}
    Every fusion-tree product-state tensor is injective at one site:
    its auxiliary matrix algebra is one-dimensional and one local
    matrix is nonzero.
\end{theorem}

\begin{proof}
    \leanok
    The auxiliary space is one-dimensional, so its matrix algebra is
    $\C$. The local matrix indexed by $h_0$ is the scalar $1$.
\end{proof}

\begin{theorem}[The action tensor is diagonal]%
    \label{thm:asymex_ising_action_diagonal}
    \lean{IsingTwist.sectorActionZ, IsingTwist.sectorAction_eq}
    \leanok
    \uses{def:asymex_ising_action, cor:asym_action_tensors}
    The action tensor of $\Theta_3$ on $\ket{\sigma1\sigma}^{\otimes L}$
    has exactly three nonzero one-site matrices:
    $2E_{88}$ at $(1,1,1)$, $2E_{99}$ at $(\psi,1,\psi)$, and
    $5E_{22}+3E_{55}$ at $(\sigma,1,\sigma)$.
    Here $E_{ij}$ denotes the matrix unit in $\MN{10}$, with
    indices starting at zero.
\end{theorem}

\begin{proof}
    \leanok
    Expand the local action of $\Theta_3$ against the product-state
    tensor at $(\sigma,1,\sigma)$. The result follows by finite
    verification of the entries in
    Definitions~\ref{def:asymex_ising_tensors}
    and~\ref{def:asymex_ising_action}.
\end{proof}

\begin{theorem}[Compression of the weighted Verlinde action]%
    \label{thm:asymex_ising_action_compression}
    % Principal declaration: IsingTwist.sectorCompression.
    \lean{IsingTwist.sectorSlots, IsingTwist.sectorLabel, IsingTwist.sectorWeights,
        IsingTwist.sectorTargets, IsingTwist.sectorGauge, IsingTwist.sectorCompression,
        IsingTwist.sectorAction_z_eq, IsingTwist.sectorAction_dim_eq}
    \leanok
    \uses{thm:asymex_ising_action_diagonal, thm:asym_multiblock}
    The action tensor of Theorem~\ref{thm:asymex_ising_action_diagonal}
    has the block triangular form of Theorem~\ref{thm:asym_multiblock}
    with four weighted bond-dimension-one product-state tensors.
    Two generate $\ket{\sigma1\sigma}^{\otimes L}$, with weights
    $5$ and $3$, from the summands indexed by $1$ and $\psi$.
    The other two generate $\ket{111}^{\otimes L}$ and
    $\ket{\psi1\psi}^{\otimes L}$, each with weight $2$, from the
    summand indexed by $\sigma$.
    The zero block has dimension six, and $10=1+1+1+1+6$.
\end{theorem}

\begin{proof}
    \leanok
    \uses{thm:asymex_explicit_gauge}
    The gauge permutes the coordinates so that $e_2,e_5,e_8,e_9$
    come first, followed by the six coordinates on which every local
    matrix vanishes. In these coordinates the action tensor is block
    diagonal with the four weighted product-state tensors and a
    six-dimensional zero block, by finite verification over integer
    entries. Block diagonality gives triangularity and identifies the
    nonzero blocks. The dimension count is clause (vii) of
    Theorem~\ref{thm:asym_multiblock}.
\end{proof}

\begin{theorem}[The extension of the action tensor splits]%
    \label{thm:asymex_ising_action_split}
    \lean{IsingTwist.sectorAction_remainder}
    \leanok
    \uses{thm:asymex_ising_action_compression}
    The remainder of the compression of
    Theorem~\ref{thm:asymex_ising_action_compression} vanishes identically.
\end{theorem}

\begin{proof}
    \leanok
    The conjugated action tensor is already block diagonal, so its
    off-diagonal blocks, which form the conjugated remainder, vanish.
\end{proof}

\begin{theorem}[The weighted Verlinde action]%
    \label{thm:asymex_ising_action_trace}
    \lean{IsingTwist.sectorAction_trace_evalWord, IsingTwist.thetaThree_mpo_sectorState}
    \leanok
    \uses{thm:asymex_ising_action_compression, prop:asym_compression_trace}
    For every nonempty word $w$ over the ten fusion-tree labels,
    \begin{align}
        \tr((\Theta_3\cdot\psi_{\sigma1\sigma})^w)
         & =(5^{|w|}+3^{|w|})\tr((\psi_{\sigma1\sigma})^w) \notag\\
         &\quad +2^{|w|}\left(\tr((\psi_{111})^w)+\tr((\psi_{\psi1\psi})^w)\right),
        \notag
    \end{align}
    where $\psi_{h_0}$ is the product-state tensor at $h_0$ and
    $\Theta_3\cdot\psi_{\sigma1\sigma}$ is the action tensor of
    Theorem~\ref{thm:asymex_ising_action_diagonal}.
    Equivalently, at every positive length $L$,
    \begin{align}
        O_L(\Theta_3)\ket{\sigma1\sigma}^{\otimes L}
         & =(5^L+3^L)\ket{\sigma1\sigma}^{\otimes L}
        +2^L\left(\ket{111}^{\otimes L}+\ket{\psi1\psi}^{\otimes L}\right).
        \label{eq:asymex_ising_verlinde}
    \end{align}
    The two contributions to $\ket{\sigma1\sigma}^{\otimes L}$ come
    from the summands indexed by $1$ and $\psi$, with different weights.
\end{theorem}

\begin{proof}
    \leanok
    Proposition~\ref{prop:asym_compression_trace} applied to the
    compression in Theorem~\ref{thm:asymex_ising_action_compression}
    gives the sum of the four weighted block traces.
    Pulling each weight out of the matrix product gives the stated
    powers. Reading this computation at each physical configuration
    gives the periodic-vector identity~(\ref{eq:asymex_ising_verlinde}).
\end{proof}
